%% The line mask, reported in two layers: the frequency coincidence, which is
%% a velocity criterion and nothing else, and the one species we do not read
%% as emission, which is an argument and is set out as one.
%% Owner: r14-mask.  Carries \label{app:linemask} (cited by 04_method.tex and
%% app_C_spatialnorm.tex) and \label{tab:masklines}.
\section{The line mask}
\label{app:linemask}

A crossing that falls on a known molecular transition has an astrophysical
explanation to hand, and the mask exists to say which crossings those are. It
reports a frequency coincidence first and an astrophysical judgement about
the coinciding species separately, so that the reproducible part can be
recomputed and the arguable part argued with.

\emph{The list, and when it was made.} The mask holds every transition of the
\FrNSpec{} species these spectral setups cover --- the three CO isotopologues,
HCN, HCO$^{+}$, CS, CN, SiO, H$_2$CO, SO, [C\,\textsc{i}] and hydrogen
recombination --- that falls inside one of the \McNIsland{} searched frequency
intervals with an upper-state energy below \McEuMax\,K. That rule is the whole
of it, and it admits \FrNTrans{} transitions (Table~\ref{tab:masklines}).
No strength cut is applied across species, since any threshold admitting the
strong lines of HCN would discard the brightest lines in these fields; within
a species the main hyperfine components are kept and the weak satellites
dropped. The list is a single Splatalogue query of CDMS and
JPL \citep{Pickett1998,CDMS2005} made on \MkListDate{}, after the disposition
criteria were fixed on \MkCritDate, and frozen before any crossing was
compared with it.

\emph{The frame is the star's own, and the sign is not the usual one.} A
crossing's frequency is carried from the topocentric value to the barycentre
using its own block's pointing and mid-time, and then to the stellar rest
frame using the star's systemic velocity. The barycentric term alone spans
\FrBarySwing\,km\,s$^{-1}$ across the survey, so a tolerance applied in the
observed frame would spend most of its budget on the Earth's motion. The
\FrEpochNEpoch{} \FrEpochStar{} epochs show
the transform working in the only way that can be checked:
\FrEpochTopoKms\,km\,s$^{-1}$ apart as observed, they agree to
\FrEpochStelKms\,km\,s$^{-1}$ in the star's frame. Toward \MkCancStar{} the
barycentric and systemic terms nearly cancel, \MkCancVBary{} against
\MkCancVSys\,km\,s$^{-1}$, so that row's two frequencies agree to
\MkCancNetkHz\,kHz although both transforms were applied.
Throughout, $\Delta v_\star = c\,(\nu_\star-\nu_0)/\nu_0$, so a positive
offset means the crossing lies \emph{above} the transition in frequency ---
the opposite sign to the velocity convention of a spectral line.

\emph{Interstellar gas is at rest in a different frame, so the mask is
evaluated in two.} Circumstellar gas shares its star's motion, but gas along
the line of sight is at rest in the local standard of rest, so the
barycentric frequency is carried to that frame as well --- one further step,
the projection of the standard solar motion --- and a crossing is attributed
if it falls inside the mask in either frame. The two offsets differ by the
star's own LSR velocity and by nothing else, which here runs
\LsVLsrLo--\LsVLsrHi\,km\,s$^{-1}$ about a median of \LsVLsrMed, the two
published conventions of the LSR differing by \LsConvKms\,km\,s$^{-1}$.
In this sample \LsNMoveWord{} crossing changes side, and none could have: the unattributed
crossing nearest a transition lies \LsNearestUnattr\,km\,s$^{-1}$ from one,
\LsGapKms\,km\,s$^{-1}$ beyond the mask and twice the largest of those
velocities. In the other direction \LsNLeaveWord{} coincidence,
\LsLeaveStar's \LsLeaveDvStar\,km\,s$^{-1}$ from \LsLeaveLine, falls outside
the mask in the LSR frame at \LsLeaveDvLsr, which is why the frames are
taken together rather than exchanged.

\emph{What the second frame changes is which attributions are interstellar.}
Of the \LsNLos{} stars carrying crossings, the composite CO($1\to0$) survey
of \citet{DameThaddeus2022}, which contains \citet{Dame2001}, covers
\LsNLosCov{} and holds emission toward \LsNLosDet{} of them, between
\LsCoVLo{} and \LsCoVHi\,km\,s$^{-1}$; the all-sky Planck CO map
\citep{PlanckCO2016} stays below \LsPlanckMax\,K\,km\,s$^{-1}$ toward every
other, including the \LsNLosOut{} outside the CfA footprint. Toward
\LsCloudStars{} the CO spectrum is non-zero at the crossing's own LSR
velocity, within \LsCloudDvMax\,km\,s$^{-1}$ of its own peak and at
\LsCloudTMin\,K or brighter against a \LsCoRms\,K channel: \LsNCloud{} of the
\MkNCoinc{} coincidences are interstellar rather than circumstellar. The
clearest is \BkAlmaName, toward the Lupus clouds: its CO($2\to1$) crossing
lies \LsLupDv\,km\,s$^{-1}$ from a \LsLupWco\,K\,km\,s$^{-1}$ cloud line at
\LsLupVco\,km\,s$^{-1}$, and the star is at \BkAlmaDistPc\,pc with the clouds
far behind it, so the field is a molecular sightline and not a disc. It is
the velocity coincidence that says so and not the field: the brightest
control position does stand above the star in all \LsLupNRingLeadWord{} of
its windows, \LsLupTCtrl{} against \LsLupTStar{} in the CO($2\to1$) one, but
so it does in \LsRingLeadA{} of the \NWinA{} Class~A windows of the survey,
and the comparison therefore distinguishes nothing. Its other \LsLupNFar{}
crossings lie at least \LsLupFarMin\,km\,s$^{-1}$ from every transition of
the list in both frames, and the mask disposes of none of them.


\begin{table}
\centering
\footnotesize
\caption{The line mask, by species: the number of transitions $N$ it holds,
the frequency range they span inside the searched intervals and their
upper-state energies. Hydrogen recombination lines carry no tabulated
$E_{\rm u}$. The criterion and the dated query above define the list
exactly, so the individual transitions can be recovered from the line
catalogues themselves.}
\label{tab:masklines}
%% GENERATED by maskso_v414.py -- do not hand-edit.
\begin{tabular}{@{}lrr@{\,}c@{\,}rr@{}}
\hline
Species & $N$ & \multicolumn{3}{c}{$\nu_0$ (GHz)} & $E_{\rm u}$ (K) \\
\hline
H & 5 & 92.03 & -- & 231.90 & -- \\
SO & 15 & 100.03 & -- & 682.68 & 16--143 \\
H$_2$CO & 13 & 101.33 & -- & 491.97 & 10--106 \\
C$^{18}$O & 3 & 109.78 & -- & 329.33 & 5--32 \\
CN & 50 & 113.12 & -- & 680.30 & 5--114 \\
CO & 3 & 115.27 & -- & 345.80 & 6--33 \\
CS & 3 & 195.95 & -- & 342.88 & 24--66 \\
SiO & 2 & 217.10 & -- & 347.33 & 31--75 \\
$^{13}$CO & 8 & 220.40 & -- & 330.59 & 16--32 \\
HCN & 7 & 354.50 & -- & 354.51 & 43--43 \\
HCO$^{+}$ & 1 & 356.73 & -- & 356.73 & 43 \\
{[C\,\textsc{i}]} & 1 & 492.16 & -- & 492.16 & 24 \\
\hline
\end{tabular}%

\end{table}

\emph{The coincidences, and the one species we do not read as emission.} Of
the \NCross{} crossings, \MkNCoinc{} fall within
$\pm\MaskHalfKms$\,km\,s$^{-1}$ of a transition of the list in their own
star's frame; that count is the first layer and all the velocity criterion
decides. Coincidences with sulphur monoxide account for \MkNSODecide{} ---
\MkSOStars, at \MkSODvList\,km\,s$^{-1}$ from transitions of upper-state
energy \MkSOEuLo--\MkSOEuHi\,K, each inside a Band~7 window that also covers
CO($3\to2$) --- and two things argue against reading those as emission.
Sulphur monoxide has not been reported in a debris disc, where the
second-generation gas is CO, C and O; and toward \FrCpStar{} the CO($3\to2$)
line falls in the same window and is absent to \FrCpCoLimMJy\,mJy in a
\FrCpCoChanKms\,km\,s$^{-1}$ channel while the crossing itself is
\FrCpFluxMJy\,mJy, so an SO line there would have to outshine an undetected
CO($3\to2$) line in the same data. Both arguments are about the species and
neither about any one crossing, so all \MkNSODecide{} are treated alike: the
offset is printed beside each in Table~\ref{tab:ledger}, where a dagger marks
the rows this applies to, and disposes of none of them, giving
\MkAttr{} attributed against \MkUnattr{} unattributed, of
which \MkRankLeadWord{} outrank every control position in their own windows.
Reading them as attributions instead gives \MkAttrAlt{} and \MkUnattrAlt,
and \MkRankLeadAlt{} such crossing; the recurrence test excludes a persisting
carrier for all \MkNSODecide{} either way, at
\MkAttrExclLo--\MkAttrExclHi$\sigma$, so the choice moves counts and no
conclusion. The reserved hold-out holds \HoNCoinc{} coincidences,
\HoNNoted{} of them the same species and treated the same way, leaving
\HoNAttr{} attributed and \HoNUnattr{} unattributed there
(\S\ref{sec:holdoutsearch}).

\emph{The half-width, and what it costs.} A one-channel test is too tight for
circumstellar gas, whose emission is shifted by many channels: the width
follows from Keplerian motion where these molecules survive,
$\lesssim$13\,km\,s$^{-1}$ beyond 10\,au, and is
$\pm\MaskHalfKms$\,km\,s$^{-1}$. At that width the mask removes
\MkMaskAGHz\,GHz of the \CvUnionA\,GHz Class~A union, \MkMaskAPct{} per cent
of it, against \FrMaskLostGHz\,GHz of the \FrMaskUnionGHz\,GHz both classes
cover, \FrMaskPct{} per cent: the cost falls on the primary experiment,
because its windows are the line-tuned ones. Evaluating the mask in two
frames widens the masked band to the union of two tubes separated by the
star's LSR velocity and raises the Class~A cost to \LsMaskTwoGHz\,GHz,
\LsMaskTwoPct{} per cent --- \LsMaskAddPct{} per cent more bandwidth, which
buys the \LsNCloud{} interstellar identifications above and no further
attribution. Between $\pm\FrLadderLo$ and
$\pm\FrLadderHi$\,km\,s$^{-1}$ the mask labels \FrAttrWTwenty{} to
\FrAttrWHundred{} of the crossings, and the \FrNRankPass{} that outrank all
\NCtrl{} control positions of their own windows are the same two at every
width.

\emph{What a bright line costs is its channels, and nothing beyond them.}
Because every above-trigger cell of a window is recorded
(\S\ref{sec:technosearch}), a weaker carrier beside a molecular line is
recorded too, and the price of line confusion is the masked channels and not
the windows holding them. Only the \MkNWholly{} fine windows narrower than
the mask itself are lost outright, \MkWhollyGHz\,GHz, \MkWhollyPct{} per
cent of the Class~A union. The one masked window that still holds a
trigger-level cell outside the mask is \MkSurvStar{} at
\MkSurvFreq\,GHz, \MkSurvDv\,km\,s$^{-1}$ from the nearest transition, and it
does not recur: the \MkSurvNRep{} other blocks covering that cell read
$T_\star=\MkSurvTRepLo$ to \MkSurvTRepHi.

\emph{The labelled crossings are tested for recurrence too.} Of the
\MkAttr{} attributed crossings, \MkAttrTested{} have a later block covering
the predicted cell and \MkAttrRecur{} are present again, toward
\MkAttrRecurStars{} at $T_\star=\MkAttrRecurTLo$ to \MkAttrRecurTHi{} --- the
test working in the positive direction on emission that is certainly real,
and no exclusion is formed for those rows, because what it would exclude has
just been observed. For the \MkAttrExcl{} not recovered a repeat window's own
largest statistic bounds the predicted cell from above, and
Table~\ref{tab:ledger} prints the row as a bound; \MkAttrUntest{} have no
covering block at all (\MkAttrUntestNames).

\emph{Image-sideband leakage accounts for no crossing.} The mask holds
signal-sideband frequencies, and the \MkImgN{} unattributed crossings that
share a block with a CO($2\to1$) crossing in the opposite sideband would need
a first local oscillator outside the range their own block's spectral windows
permit, \MkImgLOLo--\MkImgLOHi\,GHz; their image frequencies fall at least
\MkImgDvMin\,km\,s$^{-1}$ from every masked transition.

\emph{A longer list is not a better one.} Every Splatalogue transition over
these intervals with an Einstein coefficient above $10^{-5}$\,s$^{-1}$ ---
\MaskDbTrans{} transitions of \MaskDbCarriers{} species --- would mask
\MaskDbPctUnion{} per cent of the union and move \FrFullNFlip{} crossings
across the line, \FrFullNFlipRankWord{} of them rank-leading. At these
frequencies almost anything lies within $\pm\MaskHalfKms$\,km\,s$^{-1}$ of
something in such a list, which is why the species are restricted to what
these setups cover and why the outcome is a label rather than a verdict.
